\section{Extremizers and contact on the closed support}
\label{ms:extremizers}

We first show that $C_n$ is attained, by a compactness argument over all
of $\mathcal P$ that allows atoms at the threshold.

Use $\mathcal P$, $F_{n,P}(H)=\Pp(S_n\le H)$, $R_n(P)$ and $C_n$
as in Section~\ref{ms:introduction}. Throughout the article,
\[
 \beta=\E[X|X|],\qquad \gamma=\E X^3
\]
are distinct signed moments.

\subsection{Fixed-index attainment}

\begin{lemma}[Orientation and fixed-index attainment]
\label{ms:attainment}
For every fixed integer $n\ge1$, $0<C_n\le\sqrt n$. There exist
$P\in\mathcal P$ and $H\in\R$ such that
\[
 \frac{\sqrt n}{\rho(P)}
       \{F_{n,P}(H)-\Phi(H/\sqrt n)\}=C_n.
\]
Thus a maximum is attained as a positive discrepancy for a closed event.

\end{lemma}
\begin{proof}
Lyapunov's inequality gives $\rho\ge(\E X^2)^{3/2}=1$, so
$C_n\le\sqrt n$. For the symmetric law on $\{-1,1\}$, the choice $H=n$
gives $F_{n,P}(H)=1>\Phi(\sqrt n)$, proving positivity.

It suffices to maximize positive closed discrepancies over all laws. Indeed,
if $F_{n,P}(H)-\Phi(H/\sqrt n)<0$, reflection gives
\[
 \Pp(-S_n<-H)-\Phi(-H/\sqrt n)
       =\Phi(H/\sqrt n)-F_{n,P}(H).
\]
Closed thresholds increasing to $-H$ converge to the strict event on the
left. Continuity of $\Phi$ gives the same supremum using positive closed
discrepancies of the reflected law. Reflection preserves all three defining
conditions of $\mathcal P$.

Choose pairs $(P_j,H_j)$ with positive scores $r_j\to C_n$ and
$r_j\ge\eta>0$. Since a discrepancy is at most one,
\[
 1\le\rho_j\le\sqrt n/\eta=:M.
\]
The unit second moments make $(P_j)$ tight. By weak subsequence selection
on $\R$, and then compactness of $[1,M]$, we may assume
$P_j\Rightarrow P$ and $\rho_j\to b\in[1,M]$. The uniform tail estimates
\[
 \int_{|x|>A}x^2\,dP_j\le M/A,\qquad
 \int_{|x|>A}|x|\,dP_j\le M/A^2
\]
allow bounded continuous truncations to pass the first two moments to the
limit. Hence $\E_P X=0$ and $\E_P X^2=1$. Lower semicontinuity of the
third absolute moment gives $1\le\rho(P)\le b<\infty$.

The thresholds are bounded as well. If $H_j<0$, then
$F_{n,P_j}(H_j)\le n/H_j^2$ by Chebyshev's inequality. If $H_j>0$, then
the positive discrepancy is at most $1-\Phi(H_j/\sqrt n)\le n/H_j^2$.
As each discrepancy is at least $\eta/\sqrt n$, both cases give
$|H_j|\le n^{3/4}/\sqrt\eta$. Take a further subsequence with $H_j\to H$.

Because $n$ is fixed, weak convergence of the marginals implies weak
convergence of their $n$-fold products, and the continuous sum map gives
$P_j^{*n}\Rightarrow P^{*n}$. For each $\varepsilon>0$, eventually
$H_j\le H+\varepsilon$. The closed-set inequality for weak convergence gives
\[
 \limsup_j F_{n,P_j}(H_j)
 \le F_{n,P}(H+\varepsilon).
\]
Letting $\varepsilon\downarrow0$ uses right continuity of the closed CDF,
not continuity at $H$. It follows that
\[
 F_{n,P}(H)-\Phi(H/\sqrt n)\ge C_n b/\sqrt n.
\]
Therefore
\[
 C_n\ge R_n(P)\ge
 \frac{\sqrt n\{F_{n,P}(H)-\Phi(H/\sqrt n)\}}{\rho(P)}
 \ge C_n\frac b{\rho(P)}\ge C_n.
\]
All inequalities are equalities, proving attainment and $\rho(P)=b$.
In particular, extremality forces $\rho(P)=b$, so the third moments
converge without assuming uniform integrability of their third powers.
\end{proof}

\subsection{Contact on the closed support}

We now fix one attaining pair $(P,H)$. The contact inequality comes from
contaminating the raw law at a point $y$ and then standardizing it. Keeping
the raw event fixed is essential: differentiating a moving distribution
function through an atom would not be justified.

\begin{lemma}[Atom-safe contact]
\label{ms:contact}
Let $n\ge2$, $C=C_n$, $F=F_{n,P}(H)$ and $z=H/\sqrt n$ at an attaining
pair. Define $I:\R\to\R$ by
\begin{align}
 \rho I(y)={}&n^{3/2}\{F_{n-1,P}(H-y)-F\}+ny\phi(z)
   +\frac{H\phi(z)}2(y^2-1)\notag\\
 &+\frac32C\rho(y^2-1)
   -C\{|y|^3-\rho-3y\beta\}.\label{ms:contact-formula}
\end{align}
Then $I(y)\le0$ for every real $y$, and $I(y)=0$ for every point of the
closed support of $P$. In particular, that support is compact.
\end{lemma}
\begin{proof}
For fixed $y$, let $P_\varepsilon=(1-\varepsilon)P+\varepsilon\delta_y$.
Its mean and variance are
\[
 m_\varepsilon=\varepsilon y,\qquad
 v_\varepsilon=1+\varepsilon(y^2-1)-\varepsilon^2y^2.
\]
For all sufficiently small $\varepsilon\ge0$, $v_\varepsilon>0$ and the
standardized law remains in $\mathcal P$. Put
$r_\varepsilon=\E_{P_\varepsilon}|X-m_\varepsilon|^3$. The standardized
positive score evaluated at its corresponding raw threshold is
\[
 G(\varepsilon)=\frac{\sqrt n\,v_\varepsilon^{3/2}}{r_\varepsilon}
 \left\{F_{n,P_\varepsilon}(H)
  -\Phi\!\left(\frac{H-nm_\varepsilon}{\sqrt{nv_\varepsilon}}\right)\right\}.
\]
It satisfies $G(\varepsilon)\le C$ and $G(0)=C$.

The finite mixture expansion makes $F_{n,P_\varepsilon}(H)$ a polynomial
in $\varepsilon$, with derivative
$n\{F_{n-1,P}(H-y)-F\}$ at zero. No density of the sum is needed.
For the absolute third moment,
\[
 r'_0=|y|^3-\rho-3y\beta.
\]
To justify differentiation, the derivative of $|x-\varepsilon y|^3$ is
bounded, for $0\le\varepsilon\le1$, by
$3|y|(|x|+|y|)^2$, which is integrable. The contaminating term contributes
$|y|^3$ at first order. Finally the derivative of the normal argument is
\[
 -\sqrt n\,y-\frac{H}{2\sqrt n}(y^2-1).
\]
Differentiating $G$ now gives exactly $G'(0+)=I(y)$ in
\eqref{ms:contact-formula}. Maximality implies $I(y)\le0$.

Integrating against $P$ gives $\int I\,dP=0$: the convolution term
integrates to zero, as do all the displayed centered moment terms. This
integral is finite because $\rho<\infty$. The function
$y\mapsto F_{n-1,P}(H-y)$ is upper semicontinuous: at an atom its value
is the larger, closed-event value. Hence $I$ is upper semicontinuous.
If $I(y_0)<0$ at a support point, upper semicontinuity would give a
neighborhood of positive $P$-measure on which $I$ is bounded above by a
negative constant. Together with $I\le0$ elsewhere this contradicts
$\int I\,dP=0$. Thus contact holds on the entire closed support, including
atoms and accumulation points.

At this fixed pair, the right side of \eqref{ms:contact-formula} is bounded
above by $-C|y|^3+M_2y^2+M_1|y|+M_0$ for finite nonnegative constants
$M_i$. Since $C>0$, this expression is negative on both sufficiently
distant rays. Those rays contain no contact point and hence no support
point. This proves compactness.
The radius obtained here need not be uniform in $n$ or in the attaining pair.
\end{proof}

For later integral identities it is useful to record the same statement as
a CDF majorant. With $\nu$ the law of $H-S_{n-1}$, define
\begin{align}
 nQ_H(y)={}&nF+\tfrac12\{z\phi(z)+C\rho/\sqrt n\}
   -(\sqrt n\phi(z)+3C\beta/\sqrt n)y\notag\\
 &-\tfrac12\{z\phi(z)+3C\rho/\sqrt n\}y^2
   +C|y|^3/\sqrt n.\label{ms:contact-majorant}
\end{align}
Rearranging \eqref{ms:contact-formula} gives
\[
 \nu([y,\infty))\le Q_H(y),
 \qquad \nu([y,\infty))=Q_H(y)\quad
                 (y\in\operatorname{supp}P).
\]
